\section{The simulator's plant}
\label{sec:s2r}
\label{sec:s2r-plant}

The horizontal plant describes surge, sway and yaw at the center of gravity.
Its continuous dynamics and the force passed to Unity are distinguished below.

\subsection{State and frames}

The body frame is FLU at the center of gravity (CG). The body velocity is
\[
  \nuv = \begin{bmatrix} u & v & r\end{bmatrix}^\top
       \equiv \begin{bmatrix} v_x & v_y & \omega_z\end{bmatrix}^\top ,
\]
surge, sway and yaw rate: the frames section's $\bm v$ under the names used in
ship hydrodynamics. The world frame is ENU with heading
$\psi$ counter-clockwise from east, and
$\Rot(\psi) = \left[\begin{smallmatrix}\cos\psi & -\sin\psi\\ \sin\psi & \cos\psi\end{smallmatrix}\right]$
turns a body vector into the world.

\paragraph{Water-relative velocity.} Drag and added mass act on the motion
\emph{through the water}. With a constant, uniform world current $\bm c$ (ENU),
\[
  \nur = \begin{bmatrix} u_r\\ v_r\\ r\end{bmatrix}
       = \nuv - \begin{bmatrix} \Rot(\psi)^\top \bm c\\ 0\end{bmatrix}.
\]
The current does not turn the boat, so $r$ is unchanged.

\subsection{The rigid body in a turning frame}
\label{sec:s2r-rigid}

Newton's law holds in the world frame: $m\,\tfrac{d}{dt}\bm V = \bm F$, with
$\bm V = \Rot(\psi)\,(u, v)^\top$ the world velocity. The body frame turns, so
differentiating brings in the rate of turn. With $\dot\psi = r$ and
$\dot \Rot = r\,\Rot S$, $S = \left[\begin{smallmatrix}0&-1\\1&0\end{smallmatrix}\right]$,
\[
  \frac{d}{dt}\bm V = \Rot\left(\begin{bmatrix}\dot u\\ \dot v\end{bmatrix}
                     + r S \begin{bmatrix}u\\ v\end{bmatrix}\right)
                  = \Rot\begin{bmatrix}\dot u - r v\\ \dot v + r u\end{bmatrix}.
\]
In body coordinates, then,
\begin{equation}
  m(\dot u - v r) = X,\qquad m(\dot v + u r) = Y,\qquad I_z \dot r = N,
  \label{eq:s2r-rigid}
\end{equation}
where $(X, Y, N)$ is the total force and moment in the body frame. The extra
terms $-vr$ and $+ur$ are not forces; they are the price of describing a
straight-line motion in a turning frame. A boat sliding sideways at $v$ while
turning sees its ``forward'' direction swing into its velocity.

\subsection{Added mass and the Munk moment}

To accelerate a hull, the thrusters must also accelerate some of the water
around it. That water has kinetic energy, and for a hull symmetric port to
starboard and with zero sway--yaw added-mass coupling ($a_{26}=0$), it is a quadratic form in the velocity,
\[
  T_A = \tfrac12\left(a_u u_r^2 + a_v v_r^2 + a_r r^2\right),\qquad a_u, a_v, a_r \ge 0 .
\]
The $a$'s are the \emph{added masses}: the masses of water that ``come along''
in each direction. Sideways a thin hull pushes a lot of water ($a_v$ large);
end-on it slips through ($a_u$ small).

\paragraph{The lumped inertia.} The water's energy adds to the body's, so the
total is $\tfrac12\nur^\top\Min\nur$ with
\[
  \Min = \operatorname{diag}(\Min_{11}, \Min_{22}, \Min_{33}),\qquad
  \Min_{11} = m + a_u,\quad \Min_{22} = m + a_v,\quad \Min_{33} = I_z + a_r .
\]
\paragraph{Equations of motion from the energy.} The idea behind Kirchhoff's
equations is Newton's second law applied to the momentum of the body \emph{plus}
the water it carries along. That momentum is the gradient of the total kinetic
energy $T$ with respect to the velocity (for $T = \tfrac12 mv^2$ it is $mv$), so with
added mass it is $(\Min_{11}u_r, \Min_{22}v_r)$, not $m$ times the velocity. And
because the equations are written in the body frame, which turns at rate $r$, the
time derivative of any vector expressed in it picks up the rotation term
$\bm\omega\times(\cdot)$, the same correction that turns $\dot u$ into $\dot u - vr$
for the rigid body. For a body in an ideal fluid this gives the motion from the
kinetic energy $T$ alone (here $T$ is the kinetic energy, not a thrust). With
$\bm p = \partial T/\partial(u_r, v_r) = (\Min_{11}u_r, \Min_{22}v_r)$ the momentum
and $h = \partial T/\partial r = \Min_{33} r$ the angular momentum, they read
\[
  \dot{\bm p} + r S\bm p = (X, Y)^\top,\qquad
  \dot h + (u_r, v_r)\times\bm p = N .
\]
The first is the turning-frame correction of Section~\ref{sec:s2r-rigid} applied
to the momentum instead of the velocity. The second has a new term, the 2-D cross
product of velocity and momentum:
\[
  (u_r, v_r)\times(\Min_{11}u_r, \Min_{22}v_r) = u_r\Min_{22}v_r - v_r\Min_{11}u_r
  = (\Min_{22} - \Min_{11})\,u_r v_r .
\]
This is the \emph{Munk moment}. Figure~\ref{fig:s2r-munk} shows where it comes
from: when the hull moves at an angle, $\Min_{22} > \Min_{11}$ makes the momentum
point more sideways than the velocity, and a momentum that is not parallel to the
velocity carries an angular momentum about the moving centre that the hull must
balance with a yawing moment, turning it broadside to the flow.

\begin{figure}[ht]
\centering
\begin{tikzpicture}[>=Stealth, font=\footnotesize]
  \begin{scope}[rotate=20]
    \draw[fill=gray!15, rounded corners=6pt] (-1.8,-0.32) rectangle (1.8,0.32);
    \draw[dashed, gray] (-2.6,0) node[left, black] {hull axis $x$} -- (2.5,0);
    \draw[->, very thick, sA] (0,0) -- (2.2,0.62) node[below right] {velocity $(u_r, v_r)$};
    \draw[->, very thick, sB] (0,0) -- (1.8,1.3) node[above] {momentum $(\Min_{11}u_r, \Min_{22}v_r)$};
  \end{scope}
  \draw[->, thick, sC] (4.0,1.0) arc[start angle=40, end angle=-40, radius=1.3]
    node[below right, align=left] {Munk moment (clockwise here)\\ $-(\Min_{22}-\Min_{11})u_rv_r$};
\end{tikzpicture}
\caption{The Munk moment. The water is heavier to push sideways than lengthwise
($\Min_{22} > \Min_{11}$), so at an angle of attack the momentum of hull plus water
(orange) points further across the hull than the velocity (blue). The two are not
parallel, and the resulting moment turns the hull away from the flow: it is
destabilizing, and the hull's yaw damping has to beat it
in a powered hold.}
\label{fig:s2r-munk}
\end{figure}

Written out,
\begin{equation}
\begin{aligned}
  \Min_{11}\dot u_r - \Min_{22}\, v_r r &= X,\\
  \Min_{22}\dot v_r + \Min_{11}\, u_r r &= Y,\\
  \Min_{33}\dot r + (\Min_{22} - \Min_{11})\,u_r v_r &= N,
\end{aligned}
\qquad\text{i.e.}\qquad
\Min\dot\nuv_r + \Cor(\nur)\,\nur = \bm\tau,
\label{eq:s2r-fossen}
\end{equation}
where $\bm\tau = (X, Y, N)$ holds every other force (drag, thrust), and
\[
  \Cor(\nur) = \begin{bmatrix}
    0 & 0 & -\Min_{22} v_r\\
    0 & 0 & \Min_{11} u_r\\
    \Min_{22} v_r & -\Min_{11} u_r & 0
  \end{bmatrix}.
\]
If there were no water ($a = 0$), $\Min_{11} = \Min_{22} = m$ and this is exactly
\eqref{eq:s2r-rigid}: the cross term vanishes because a rigid body has the same
mass in every direction.

\paragraph{The Munk moment.} The yaw equation has the moment
$-(\Min_{22} - \Min_{11})\,u_r v_r = -(a_v - a_u)\,u_r v_r$ with no external cause.
For a boat going forward ($u_r > 0$) with a little leeway to port ($v_r > 0$) it is
negative: it turns the bow to starboard, away from the flow, which makes the leeway
angle larger. Any elongated body in a flow feels it; it is why a hull with too
little keel or skeg is directionally unstable, and why the yaw damping of the hull
matters as much as its drag.

\paragraph{A check: $\Cor$ does no work.} The power of the $\Cor$ terms is
$\nur^\top\Cor(\nur)\nur
= -\Min_{22}u_r v_r r + \Min_{11} u_r v_r r + \Min_{22} u_r v_r r - \Min_{11} u_r v_r r = 0$.
They only redirect momentum, as they must: water without friction cannot drain
energy. The discrete step is checked for non-growing energy; it does not exactly
conserve the continuous kinetic energy.

\paragraph{From ground to water velocity, exactly.} The rigid-body part can be
written with either velocity, because for a constant current
\[
  m(\dot u - v r) = m(\dot u_r - v_r r),\qquad m(\dot v + u r) = m(\dot v_r + u_r r).
\]
Proof: the current in body coordinates is $\nuv_c = \Rot(\psi)^\top\bm c$, and
$\tfrac{d}{dt}\Rot^\top = -r S\Rot^\top$, so $\dot u_c = r v_c$ and
$\dot v_c = -r u_c$. Then
$\dot u_r - v_r r = \dot u - r v_c - (v - v_c)r = \dot u - v r$, and likewise for
sway. A simulator can therefore integrate the rigid body on ground velocity, as a
physics engine does, and apply added mass, Coriolis and damping on $\nur$
with the force mapping below.

\paragraph{Sway--yaw coupling.} The implemented model also supports $a_{26}$.
Writing $a_u=a_{11}$, $a_v=a_{22}$ and $a_r=a_{66}$, its mass and added
Coriolis terms are
\[
 M_A=\begin{bmatrix}a_u&0&0\\0&a_v&a_{26}\\0&a_{26}&a_r\end{bmatrix},
 \quad M_{RB}=\operatorname{diag}(m,m,I_z),\quad \Min=M_{RB}+M_A,
 \quad C_A(\nur)\nur=\begin{bmatrix}
 -(a_vv_r+a_{26}r)r\\a_uu_rr\\(a_v-a_u)u_rv_r+a_{26}u_rr
 \end{bmatrix}.
\]
The continuous ground-velocity equation is
$\Min\dot\nuv=\bm\tau+M_A\dot\nuv_c-C_A(\nur)\nur-\bm D(\nur)-C_{RB}(\nuv)\nuv$,
where $\nuv_c=(R^\top\bm c,0)$ and
$C_{RB}(\nuv)\nuv=(-mvr,mur,0)^\top$.
The sway--yaw solve uses the coupled $2\times2$ block. The displayed diagonal
Kirchhoff equations specialize this model to $a_{26}=0$.

\subsection{Hull damping}

Each horizontal axis $s\in\{u_r,v_r,r\}$ has sign-specific nonnegative linear
and quadratic coefficients. The hull feels the negative of
\[
 D_j(s)=(D_{1,j}^{\operatorname{sgn}s}+D_{2,j}^{\operatorname{sgn}s}|s|)s.
\]
Heave, roll and pitch use symmetric linear-plus-quadratic body-frame damping.

The damping derivative used by the implicit mapping takes the stronger linear
coefficient at zero:
\[
 J_j(s)=\begin{cases}D_{1,j}^{\operatorname{sgn}s}+2D_{2,j}^{\operatorname{sgn}s}|s|,&s\ne0,\\
 \max(D_{1,j}^+,D_{1,j}^-),&s=0.\end{cases}
\]

\subsection{The force applied by Unity}

For step $\Delta t$, the damping correction uses the diagonal effective inertias
$M_{\mathrm{eff},j}=\Min_{jj}$. Its exponential term is
\[
 x_j=\frac{\Delta t J_j}{M_{\mathrm{eff},j}},\qquad
 h_j=M_{\mathrm{eff},j}\left(\frac{x_j}{1-e^{-x_j}}-1\right),
 \qquad h_j\simeq M_{\mathrm{eff},j}(x_j/2+x_j^2/12)\quad(x_j<10^{-4}).
\]
The implementation returns zero for nonpositive $x_j$ and uses the displayed
series for positive small $x_j$.

Let $H=\operatorname{diag}(h_j)$ and collect the actuator and hydrodynamic terms
into $\bm\tau_h$. The force correction includes $H$ in both matrices:
\[
 \bm\tau_h=\bm\tau+M_A\dot\nuv_c-C_A(\nur)\nur-\bm D(\nur),\qquad
 \widetilde M=\Min+H,\qquad
 \bm F_U=\bm\tau_h-(M_A+H)\widetilde M^{-1}
          (\bm\tau_h-C_{RB}(\nuv)\nuv).
\]
Unity applies $\bm F_U$ to its rigid mass and inertia. Wind, when enabled, is a
separate force Unity applies directly at its point of application; it does not
enter $\bm\tau$ or this correction. The pure step updates
world velocity and yaw rate before position and heading (symplectic Euler).

\subsection{Thrusters}

For clamped pulse $p\in[1100,1900]$\,$\mu$s, neutral offset $o$ and deadband
$d_b$, the effort is zero when $|p-1500-o|\le d_b$ and otherwise is
$q=\operatorname{clip}((p-1500-o)/400,-1,1)$. The range outside the deadband
is not rescaled. After the command delay the static target is
\[
 f_{\mathrm{ss}}=\begin{cases}K_fq^2,&q>0,\\-K_rq^2,&q<0,\\0,&q=0,\end{cases}
 \qquad f_{\mathrm{ss}}\leftarrow f_{\mathrm{ss}}(\Vbus/V_{\mathrm{ref}})^{n_V}.
\]
Voltage scaling applies only for nonzero $n_V$ and finite positive supplied
voltage. Nonfinite pulse input is replaced with neutral.

The delay is rounded to whole ticks and the first-order lag uses its exact
sampled gain:
\[
 N_d=\max(0,\operatorname{round}(t_d/\Delta t)),\qquad
 f_{k+1}=f_k+(1-e^{-\Delta t/\tau_T})(f_{\mathrm{ss},k-N_d}-f_k).
\]
For $\tau_T\le0$ the gain is one. A step-size change resizes and clears the
delay buffer. Mount directions and lever arms turn individual thrusts into a
body force and moment.

\file{crane\_ml: Assets/Scripts/Physics/Water/Dynamics/MarinePlanarModel.cs, MarineHydrodynamics.cs; Assets/Scripts/Actuators/Motors/ThrusterModel.cs}
